% ASCII LaTeX source
% Title: McBride Derivatives and Quantale Structure in the Chaos Language Algorithm
% Author: ZeroBot / ProtomegaTron
% Date: 2026-07-05
%
% Compile: tectonic mcbride_cla_analysis.tex
%
\documentclass[11pt]{article}
\usepackage[utf8]{inputenc}
\usepackage[T1]{fontenc}
\usepackage{amsmath,amssymb,amsthm}
\usepackage{geometry}
\usepackage{hyperref}
\usepackage{enumitem}
\usepackage{booktabs}

\geometry{margin=1in}

\newtheorem{theorem}{Theorem}
\newtheorem{proposition}{Proposition}
\newtheorem{definition}{Definition}
\newtheorem{conjecture}{Conjecture}
\newtheorem{example}{Example}
\newtheorem{corollary}{Corollary}

\title{McBride Derivatives and Quantale Structure in the\\Chaos Language Algorithm}
\author{ProtomegaTron \and ZeroBot}
\date{5 July 2026}

\begin{document}
\maketitle

\begin{abstract}
We analyze the Chaos Language Algorithm (CLA) through the lens of McBride derivatives and quantale-theoretic structure. The McBride (discrete) derivative of the MDL objective with respect to candidate grammar edits provides a local sensitivity measure that enables efficient pre-ranking of proposals. Pattern intensity in a quantale of choice generalizes the binary context-matching currently used for category induction, yielding a graduated similarity measure. Emergent synergy---when chunk and category edits compose non-additively---suggests a joint-proposal extension to the greedy single-edit search. We further connect the grammar trajectory to delay-logistic dynamics predicted by the quantale ODE framing, giving convergence diagnostics. These results yield five concrete practical implications for the \texttt{chaoslang} implementation, which we enumerate.
\end{abstract}

\tableofcontents
\newpage

\section{Background}

\subsection{The Chaos Language Algorithm}

The CLA converts a symbolic stream $S$ derived from a chaotic trajectory into a grammar $G$ that permits exact reconstruction of $S$. It alternates two kinds of edits:
\begin{itemize}
    \item \textbf{Chunk edits} $(i,j)$: introduce a nonterminal $N \to S[i..j]$ and replace all occurrences of that block with $N$.
    \item \textbf{Category edits}: group distinct symbols $x, y$ into a meta-symbol $M = \{x, y\}$, recording per-occurrence member selections in a side table $M[v]$ to preserve exact reconstruction.
\end{itemize}
Acceptance is governed by Minimum Description Length (MDL): an edit is accepted iff $\Delta L < 0$, where $L(G, S)$ is the total description length of grammar plus stream encoding.

\subsection{McBride Derivatives}

Following McBride's discrete calculus, the \emph{derivative} of a function $f$ at a point $a$ with respect to an increment $\delta$ is:
\[
\partial_\delta f(a) \;=\; f(a + \delta) - f(a).
\]
In the CLA context, we take $f = L$ (the MDL objective), $a = (G, S)$ the current grammar-stream pair, and $\delta = (i,j)$ a candidate chunk edit. The McBride derivative $\partial_{(i,j)} L$ measures the marginal change in description length from excising the block $S[i..j]$.

\subsection{Quantales and Pattern Intensity}

A \emph{quantale} $(Q, \odot, \bigvee, \bot)$ is a complete lattice equipped with an associative monoidal operation $\odot$ distributing over arbitrary joins. The CLA MDL objective lives in the additive real-valued quantale $([0,\infty], \cdot, +, 0)$. An alternative is the \emph{probability quantale} $([0,1], \times, \max, 0)$.

Given a candidate category merge of symbols $y, z$ in context $m$, the \emph{pattern intensity} is:

\begin{definition}[Pattern Intensity]
\[
I_{y,z}(m) \;=\; \sigma_{y,z}(m) - e_{y,z}(m),
\]
where $\sigma_{y,z}(m)$ is the emergent synergy (non-additive compression gain) from merging $y$ and $z$ in context $m$, and $e_{y,z}(m)$ is the entropy cost of the member-selection side table.
\end{definition}

When $I_{y,z}(m) > 0$, the merge is beneficial; the magnitude gives a graduated measure rather than a binary accept/reject.

\section{McBride Derivative Analysis of CLA}

\begin{theorem}[McBride Derivative as Proposal Pre-Ranker]
\label{thm:mcbride-prerank}
Let $L(G, S)$ be the MDL objective and $(i, j)$ a candidate chunk edit. The McBride derivative
\[
\partial_{(i,j)} L \;=\; L(G + (i,j),\, S) - L(G, S)
\]
equals the exact MDL delta from accepting edit $(i,j)$. Therefore, ranking candidate edits by $\partial_{(i,j)} L$ produces the same ordering as full MDL evaluation, while permitting early pruning of candidates whose derivative exceeds zero without computing the full post-edit encoding.
\end{theorem}

\begin{proof}
By definition, $L(G + (i,j), S)$ is the description length after accepting the chunk edit. The MDL acceptance criterion is $\Delta L = L(G + (i,j), S) - L(G, S) < 0$, which is exactly $\partial_{(i,j)} L < 0$. Since the derivative \emph{is} the delta, the ranking is identical. The computational advantage comes from the fact that $\partial_{(i,j)} L$ can be approximated locally---counting block occurrences and estimating the rule cost---without fully re-encoding the stream and grammar.
\end{proof}

\begin{corollary}[Cheap Pre-Ranking]
\label{cor:prerank}
An approximate McBride derivative $\widehat{\partial}_{(i,j)} L$ computed from local occurrence counts and rule-length estimates preserves the ranking of Theorem~\ref{thm:mcbride-prerank} up to approximation error. One may compute $\widehat{\partial}$ for all candidates, retain the top-$k$ by $\widehat{\partial}$, and evaluate full MDL only on those $k$, reducing per-iteration cost from $O(n \cdot |S|)$ to $O(n + k \cdot |S|)$ where $n$ is the number of candidates.
\end{corollary}

\section{Emergent Synergy and Joint Proposals}

\begin{proposition}[Emergent Synergy as Joint Proposal Criterion]
\label{prop:synergy}
Let $y$ be a chunk edit and $z$ a category edit. Define the emergent synergy:
\[
\sigma_{y,z}(m) \;=\; \Delta L(y \circ z) - \Delta L(y) - \Delta L(z),
\]
where $\Delta L(\cdot)$ is the MDL delta of applying the edit in isolation and $y \circ z$ denotes simultaneous application. When $\sigma_{y,z}(m) > 0$, the joint application compresses more than the sum of individual applications. Greedy single-edit search misses such pairs.
\end{proposition}

\begin{example}[Synergy Necessitates Joint Proposals]
Consider the stream $S = \texttt{axb\,ayb\,axb}$. Chunking $\texttt{aMb}$ with $M = \{x, y\}$ yields compression, but only if the category $M$ is introduced first (or simultaneously). Chunking $\texttt{axb}$ alone fails because the block appears only twice (insufficient for MDL gain). Categorizing $M = \{x, y\}$ alone fails because without a chunk wrapper, the category saves no sequential structure. The synergy $\sigma_{y,z} > 0$ signals that these edits must be proposed jointly.
\end{example}

\begin{conjecture}[Synergy Detection via Rejected Proposals]
\label{conj:synergy-detect}
When a chunk proposal is rejected because the candidate block is not repeated enough, one can check whether a category edit would unify near-matching blocks. If the post-category block count reaches the MDL threshold, the pair $(\text{chunk}, \text{category})$ is
 a synergy candidate with $\sigma > 0$. This gives a constructive detection procedure implementable as a post-rejection hook.
\end{conjecture}

\section{Quantale Pattern Intensity}

\begin{definition}[Quantale Pattern Intensity]
\label{def:quantale-intensity}
In a quantale $(Q, \odot, \bigvee, \bot)$, the pattern intensity of merging symbols $y, z$ in context $m$ is:
\[
I_{y,z}^{Q}(m) \;=\; \sigma_{y,z}(m) \;\ominus_Q\; e_{y,z}(m),
\]
where $\ominus_Q$ is the residuation (implication) in $Q$. For the additive quantale, $\ominus_Q = -$ (subtraction); for the probability quantale, $\ominus_Q$ is related to the Boolean implication.
\end{definition}

This generalizes the binary context-signature match in the current \texttt{ContextCategoryInducer} to a graduated measure: pairs with $I_{y,z}(m)$ above a threshold are proposed, ranked by intensity. The Jensen-Shannon divergence threshold $\theta$ in the CLA specification already gestures at this; the quantale formalism gives it a principled algebraic form.

\begin{proposition}[Quantale Choice as Tunable Knob]
\label{prop:quantale-knob}
The additive quantale $([0,\infty], \cdot, +, 0)$ (MDL) emphasizes cumulative compression: all repeated patterns contribute. The probability quantale $([0,1], \times, \max, 0)$ emphasizes the strongest single pattern: $\bigvee = \max$ means only the largest contribution survives. For highly periodic trajectories, the additive quantale may over-fragment (many small chunks); the probability quantale would favor the dominant period. For chaotic-transient trajectories with multi-scale structure, the additive quantale captures hierarchical nesting better. Different quantales thus produce different grammars, and the choice is a domain-tunable knob.
\end{proposition}

\section{Grammar Trajectory Dynamics}

\begin{proposition}[Grammar Trajectory as Delay-Logistic Dynamics]
\label{prop:trajectory}
Let $G(t)$ be the grammar after $t$ accepted edits and $L(t) = L(G(t), S)$. The CLA greedy loop defines a discrete dynamical system on the grammar space. Because proposal statistics at iteration $t$ depend on $G(t-1)$, the system has delayed feedback. In the quantale ODE framing, this connects to the delay-logistic dynamics of WT \S22.6: the MDL score trajectory $\{L(t)\}$ can exhibit monotone decrease, oscillation, or chaotic transients depending on the acceptance threshold and the structure of the candidate space.
\end{proposition}

\begin{corollary}[MDL Trajectory as Convergence Diagnostic]
\label{cor:diagnostic}
Logging $\{L(t)\}_{t=0}^{T}$ and checking for oscillation or non-monotone behavior provides a free diagnostic: oscillation indicates the greedy search is cycling between competing grammars, which is qualitatively different from convergence. This guides tuning of acceptance thresholds, cooling schedules, or synergy-aware proposal strategies.
\end{corollary}

\section{Practical Implications for \texttt{chaoslang}}

We enumerate five concrete engineering improvements, ordered by implementation cost.

\subsection{1. Log Per-Iteration MDL Trajectory (Cost: $\sim$10 lines)}

Add a list field to the greedy loop that appends $L(t)$ after each iteration. Plot or inspect the trajectory. This gives immediate convergence diagnostics per Corollary~\ref{cor:diagnostic}. Detecting oscillation tells us the greedy search is stuck cycling between competing grammars---qualitatively different from having converged.

\textbf{Implementation:} In \texttt{scoring.py} / \texttt{CLA.simple()}, add \texttt{self.mdl\_trajectory = []} and append \texttt{current\_score.total\_bits} after each accepted or rejected iteration.

\subsection{2. Log Pattern Intensities for Accepted Categories (Cost: small)}

Alongside the existing JS-divergence computation in \texttt{ContextCategoryInducer}, compute $I_{y,z}(m)$ per Definition~\ref{def:quantale-intensity} for each candidate category pair. Compare rankings: does intensity rank candidates differently from JS-divergence? If so, intensity may be a better proposal filter.

\textbf{Implementation:} Add an \texttt{intensity} field to category proposals in \texttt{categorization.py}; compute $\sigma - e$ from occurrence counts and member-table entropy estimates.

\subsection{3. Flag Synergy Candidates (Cost: moderate)}

When a chunk proposal is rejected because the block is not repeated enough, check whether a category edit would unify near-matching blocks per Conjecture~\ref{conj:synergy-detect}. If the post-category block count reaches the MDL threshold, log the pair as a synergy candidate. This directly addresses the known weakness of greedy single-edit search getting stuck in local minima.

\textbf{Implementation:} Add a post-rejection hook in the greedy loop that queries the category proposer for potential unifying merges on the rejected chunk's near-matches.

\subsection{4. McBride Derivative Pre-Ranking (Cost: moderate refactor)}

Replace the current full-MDL evaluation for all candidates with an approximate McBride derivative $\widehat{\partial}_{(i,j)} L$ computed from local occurrence counts and rule-length estimates (Corollary~\ref{cor:prerank}). Retain top-$k$ candidates by $\widehat{\partial}$, evaluate full MDL only on those $k$. The payoff is real for scaling to OmegaSim-dimensional trajectories where the candidate space is large.

\textbf{Implementation:} Add a \texttt{approximate\_delta} method to \texttt{TwoPartMDLScorer} that estimates $\Delta L$ without full stream re-encoding. Use it in the proposal filter before the exact \texttt{delta()} call.

\subsection{5. Quantale Choice as Domain-Tunable Knob (Cost: speculative, testable)}

Swap the scoring algebra from the additive quantale to the probability quantale and compare grammar quality on logistic-map (periodic) vs.\ Lorenz (chaotic-transient) trajectories. The probability quantale emphasizes the strongest single pattern rather than cumulative compression, which may produce cleaner grammars for highly periodic streams.

\textbf{Implementation:} Abstract the scoring operation behind a \texttt{Quantale} protocol with \texttt{combine} ($\odot$) and \texttt{aggregate} ($\bigvee$) methods. Provide \texttt{AdditiveQuantale} (current) and \texttt{ProbabilityQuantale} implementations. Run comparative benchmarks.

\section{Summary and Open Questions}

\begin{itemize}
    \item The McBride derivative \emph{is} the MDL delta (Theorem~\ref{thm:mcbride-prerank}), so it can replace full evaluation for pre-ranking without loss of accuracy---only approximation error from the local estimate.
    \item Emergent synergy $\sigma > 0$ (Proposition~\ref{prop:synergy}) identifies joint chunk+category proposals that greedy single-edit search misses. Constructive detection via rejected-proposal hooks is conjectured (Conjecture~\ref{conj:synergy-detect}).
    \item Quantale pattern intensity (Definition~\ref{def:quantale-intensity}) generalizes binary context matching to a graduated measure, with the quantale choice itself being a tunable knob (Proposition~\ref{prop:quantale-knob}).
    \item Grammar trajectory dynamics (Proposition~\ref{prop:trajectory}) connect to delay-logistic dynamics, and the MDL trajectory is a free convergence diagnostic (Corollary~\ref{cor:diagnostic}).
\end{itemize}

\subsection*{Open Questions}
\begin{enumerate}
    \item What is the tightest bound on the approximation error of $\widehat{\partial}_{(i,j)} L$ relative to the exact MDL delta, as a function of stream length and grammar complexity?
    \item Can synergy detection be made complete---i.e., are there synergy pairs that the rejected-proposal hook misses?
    \item How does the optimal quantale choice depend on the Lyapunov exponent or other dynamical invariants of the source trajectory?
    \item Does the grammar trajectory exhibit actual chaotic transients (as opposed to simple oscillation) for certain stream classes, and if so, what does this imply about the existence of multiple locally optimal grammars?
    \item Can the quantale ODE framing yield a \emph{predictive} model of convergence time, rather than merely a diagnostic after the fact?
\end{enumerate}

\end{document}
